\documentclass[12pt]{article}
\usepackage[utf8]{inputenc}
\usepackage[T1]{fontenc}
\usepackage[margin=1in]{geometry}
\usepackage{setspace}
\usepackage{xcolor}
\usepackage[colorlinks=true, urlcolor=blue!60!black, linkcolor=blue!60!black]{hyperref}
\setlength{\parindent}{0em}
\setlength{\parskip}{0.6em}
\renewcommand{\baselinestretch}{1.0}
\begin{document}
\begin{center}
{\LARGE \textit{Humphrey's Executor v.\ United States}, 295 U.S.\ 602 (1935)\par}
\vspace{0.5em}
{\large Case Brief --- Two-Page Edition\par}
\end{center}

\section{Memory Jogger}
Congress could limit the President to ``for-cause'' removal of FTC Commissioners, whose duties were ``quasi-legislative'' and ``quasi-judicial''; overruled in 2026 by \textit{Trump v.\ Slaughter} (\href{https://supreme.justia.com/cases/federal/us/295/602/}{Justia}).

\section{Detailed Case Facts}
The FTC Act gave commissioners staggered seven-year terms and allowed removal ``for inefficiency, neglect of duty, or malfeasance in office.'' President Roosevelt removed Commissioner William Humphrey in 1933 over policy disagreements, citing none of those causes. Humphrey insisted he remained in office until his death.

\section{Procedural History}
Humphrey's executor sued in the Court of Claims for back salary; that court certified two questions to the Supreme Court, which answered both ``Yes.''

\section{Judicial Votes}
Unanimous. Sutherland wrote for eight Justices; McReynolds agreed in a separate statement referring to his \textit{Myers} dissent.

\section{Holding}
The FTC Act barred removal except for the listed causes, and that limit was constitutional. Whether removal may be restricted depends on ``the character of the office'': \textit{Myers} covers only purely executive officers.

\section{Analysis of Opinions}
\textbf{Majority:} \textit{Myers} involved a postmaster; its broader statements were dicta. The FTC carried out legislative policies, reported to Congress, and adjudicated, so it was no ``arm or an eye of the executive.'' An officer serving at the President's pleasure cannot act independently, and unlimited removal power would let the Executive coerce the other branches. \textbf{McReynolds:} the President has no illimitable removal power over statutory officers.

\section{Examples: Future Applications}
\textbf{Same-Side:} a multi-member claims commission performing quasi-judicial duties (\textit{Wiener}). \textbf{Opposite-Side:} a single-director agency with enforcement power (\textit{Seila Law}). \textbf{Fence-Sitter:} the Federal Reserve Board, carved out in \textit{Trump v.\ Cook}.

\section{Critique}
Unitary-executive critics argued that Article~II vests all executive power in the President and that the FTC wielded real enforcement power; that view prevailed in \textit{Slaughter}. Defenders say history and function support independent expert agencies.

\section{Key Quotations}
\begin{quote}
``[O]ne who holds his office only during the pleasure of another cannot be depended upon to maintain an attitude of independence against the latter's will.''
\end{quote}
\noindent --- Justice Sutherland, for the Court, 295 U.S.\ at 629.

\end{document}
